%% 第六章 总结与展望

\chapter{总结与展望}
\chaptermark{总结与展望}

本文围绕长期服务机器人情景记忆问答中的检索效率、存储膨胀和错误修正问题展开研究。随着机器人从短时任务执行走向长期连续部署，系统需要持续积累多模态经历，并在用户提出回溯性问题时从大量历史记忆中定位相关证据。现有长上下文语言模型、检索增强生成和层级记忆方法分别在上下文利用、语义检索和层级组织方面提供了可行思路，但仍难以同时解决长历史检索效率下降、记忆规模持续增长以及错误信息写入后难以修正等问题。针对上述问题，本文提出 GRAF-Mem，一个基于层级记忆树的主动式长期情景记忆框架，从图增强检索、效用引导遗忘和反馈驱动修正三个方面提升长期记忆系统的可扩展性与可维护性。

首先，针对层级记忆树主要依赖纵向父子关系、难以利用跨事件横向关联的问题，本文提出事件关联图增强检索方法。该方法从 L2 事件节点自动构建包含时间、物体、位置、动作和因果关系的事件关联图，并在语义种子召回的基础上引入图邻居扩展与联合重排机制，使检索过程能够沿结构化关系发现相关证据。同时，本文将任务、物体、日期和时间邻居等信息封装为结构化检索工具，降低了智能体在长历史记忆树中盲目搜索的成本。实验结果表明，在 TEACh 长历史问答任务中，GRAF-Mem 在提升语义正确率的同时显著降低了 token 消耗，说明事件关联图能够有效弥补纯文本相似度检索在横向证据发现方面的不足。

其次，针对长期运行中记忆持续追加导致的存储膨胀问题，本文提出效用引导的渐进式遗忘方法。该方法综合考虑事件的新近性、语义独特性和多信号重要性，并引入图中心性刻画事件在记忆结构中的关联价值。根据效用评分，系统对不同事件执行完整保留、去除原始多模态大对象或仅保留摘要的分级压缩策略，同时通过对话豁免、失败豁免、目标边界保护和最低保留比例等机制降低关键证据被误删的风险。实验结果表明，在保持问答结果基本稳定的前提下，渐进式遗忘能够有效减少记忆树中的冗余信息；在真实机器人数据上的补充验证进一步说明，图中心性有助于将有限的保留预算分配给更具结构价值的事件，从而实现更具选择性的记忆压缩。

再次，针对长期记忆中错误信息难以修正并可能沿层级摘要传播的问题，本文提出反馈驱动的记忆修正机制。该机制包括错误定位、LLM 辅助最小化修正、传播检测和自动传播修正四个阶段，并通过非侵入式视图重定向机制在不破坏原始记忆数据的前提下挂载修正结果。受控错误注入实验表明，该机制能够根据用户反馈定位并修正错误事实，在连续帧污染场景中检测同源错误传播，并通过端到端闭环验证将用户反馈转化为后续答案改善。上述结果说明，长期记忆系统不仅应具备"存储"和"检索"能力，还应具备面向长期运行的维护能力。

综合来看，本文的主要贡献在于将长期情景记忆从被动的历史保存与检索，推进到包含检索、压缩和修正的主动维护过程。事件关联图在三个模块中发挥了统一的信息基础作用：在检索阶段提供横向扩展路径，在遗忘阶段提供结构重要性信号，在修正阶段提供传播检测依据。通过这一统一结构，GRAF-Mem 将原本相对独立的检索效率、存储管理和错误维护问题纳入同一框架，为长期服务机器人情景记忆问答提供了一种系统化解决方案。

同时，本文工作仍存在一定局限。首先，当前事件关联图中的边类型和构建阈值主要依赖人工设定，图结构的自适应能力仍有限。其次，遗忘策略中的阈值和权重尚未根据实时存储压力、任务密度和用户查询习惯进行动态调整。再次，修正模块主要基于文本摘要和受控错误注入实验进行验证，对于真实用户反馈中的模糊表达、噪声反馈和跨时间错误传播仍需要更充分的实验支持。此外，本文实验主要基于 TEACh 和 ARMARX 数据，尚未在真实家庭环境中进行长周期在线部署验证。

未来工作可以从以下几个方向继续展开。第一，可以进一步增强事件关联图的表达能力，引入更丰富的任务关系、因果关系和用户意图关系，并探索根据验证集表现自动调整边权重和检索策略的方法。第二，可以研究自适应遗忘机制，使压缩强度能够根据存储预算、近期任务密度和历史查询分布动态变化，从而在不同部署阶段实现更灵活的存储管理。第三，可以将修正模块扩展为多模态错误定位机制，结合文本摘要、视觉特征和场景图信息共同判断错误来源，并进一步支持跨 episode 的系统性错误传播检测。第四，可以将 GRAF-Mem 集成到真实机器人软件栈中，在自然语言交互、长期运行和真实用户反馈条件下验证其稳定性、鲁棒性和工程可用性。

总之，本文针对长期服务机器人情景记忆问答中的关键问题，提出并验证了图增强检索、效用引导遗忘和反馈驱动修正三类方法。实验结果表明，主动式记忆维护有助于提升长期记忆系统的效率、经济性和可靠性。随着服务机器人逐步进入更加开放和长期的应用场景，具备可检索、可压缩、可修正能力的情景记忆系统将成为提升人机长期协作体验的重要基础。
